% ==============================================================================
% EM tenure with point-mass + Exp contamination (epsilon model)
% Created: 2026-04-30
%
% This document presents a structural reformulation of the tenure
% measurement model that replaces the Gaussian (sigma^2_g) measurement-error
% component with a point-mass + Exponential contamination mixture, and
% replaces the EMG spell distribution with a plain Exponential. The model
% also introduces a spell-pair joint emission to allow tenure flanks to
% identify misclassification.
%
% This document is a self-contained companion to "EM tenure.tex". It does
% not modify the base specification.
% ==============================================================================
\documentclass[12pt]{article}
\usepackage[margin=2cm]{geometry}
\setlength{\parindent}{1em}
\setlength{\parskip}{1em}

\usepackage{amsmath}
\usepackage{amssymb}
\usepackage{booktabs,dcolumn,setspace,float}

\begin{document}

\title{Contamination model: A point-mass + exponential contamination model for misclassified employment dynamics with tenure}
\author{}
\date{April 2026}
\maketitle

\section{Motivation}

The base EM-tenure specification (\texttt{EM tenure.tex}) relies on a
Gaussian increment density \(N(0, 2\sigma_g^2)\) for the tenure-measurement
component, an EMG \(f_{\mathrm{EMG}}(\cdot;\lambda_g,\sigma_g^2)\) for spell
levels, and a single error parameter \(\pi\).\footnote{(The
duration-contamination extension \texttt{EM tenure rho.tex}) also relies on a Gaussian increment density, and relies on $\pi$ and $\rho$ } Three empirical facts in the QLFS rule out the Gaussian
measurement model:

\begin{enumerate}
\item \textbf{Tenure is integer-month, not continuous Gaussian.} 
  Reported tenures are integer-month value which does not concur with the
  \(\mathcal{N}(0,\sigma_g^2)\) noise assumption.
\item \textbf{63.2\% of E\(\to\)E continuation increments are exactly 0.25
  years}. The Gaussian MLE for \(\sigma_g^2\)
  collapses to zero. Among the remaining 36.8\%: 1.5\% are within 0.5 months
  of 3, 7\% are within 1--6 months, and 28\% are wildly off (negatives,
  year-shifted, or large positive deviations). This is a so-called \textit{contamination
  distribution} \textemdash i.e. a mixture \textemdash not a Gaussian additive error.
\item \textbf{36.7\% of within-panel, new job entrants (state pattern
  \(s_{t-1}=0,s_t=1\)) report tenure \(>\) 12 months.} A genuine job entry
  cannot have 12+ months of tenure. These observations must be
  misclassifications at \(t-1\): the respondent was actually employed at
  \(t-1\) but misclassified, and the tenure at wave \(t\) reflects the
  ongoing spell (or the tenure is wrong).
\end{enumerate}

Fact (3) is an important justification for this \(\varepsilon\) tenure model: tenure
data corroborates employment-state Markov violations in identifying misclassification.
The base endogenous EM tenure model fails to exploit this corroboration
fully. The base EM tenure model collapses \(\sigma_g^2 \to 0\) and routes all
contamination through \(\pi\) (which inflates to over 25\%).\footnote{The \(\rho\) model
restores plausible \(\pi \approx 4\%\) but absorbs all tenure inconsistency
into \(\rho\), so tenure does not actively inform \(\pi\) beyond what
state-Markov violations alone provide.} It models the wave-3 tenure under
the latent history \(h=(1,1,1)\) with intermediate misclassification as an
EMG draw \emph{independent} of \(g_1\), discarding the structural information
that waves 1 and 3 belong to the same spell. 

\textbf{Importantly}, in the base EM tenure model,
random EMG draws for tenure can only occur when there has truly been a misclassification
of employment status. However, we note that it is possible that employment status is reliable
even when tenure is not. Moreover, even if an individual is completely mismatched such
that tenure becomes a random draw from the tenure distribution, fthere would still be an
approximately 50\% chance that their employment status is correct by coincidence due 
to the employment rate being near 50\%. Under the base EM tenure model, it is required that
for these observations true employment be different to observed employment, and to accommodate
that it drives up both mislclassification and true transitions. For example, if \(h=(1,1,1)\) 
and \(s=(1,1,1)\) but \(g=(1, 5, 1.5)\), the base EM tenure model would see that 
tenure in wave 2 is inconsistent with 0.25 year increments, and the would not explain this
tenure pattern using \(\sigma\) because it has collapsed to zero as discussed above. The only
other way to account for this tenure is to declare a misclassification at wave 2, 
which requires thinking that \(h=(1,0,1)\) and \(s=(1,1,1)\) such that there is both a transition into
and out of true employment, but that there was a misclassification at wave 2. In this way, 
both the true transitions and the misclassification estimates are driven up.

This contamination specification addresses all three issues. 
It replaces the Gaussian + EMG measurement
model with a point-mass + Exponential contamination mixture, eliminating
\(\sigma_g^2\) entirely, and introduces a \emph{spell-pair} joint emission
that uses tenure flanks to discriminate misclassification from genuine
transition.

\section{Notation and terminology}

We follow \texttt{EM tenure.tex} notation.
Waves \(t \in \{1,2,3\}\), latent state \(h_{t} \in \{0,1\}\), observed
state \(s_{t} \in \{0,1\}\), tenure \(g_{t} \ge 0\), nonemployment
duration \(d_{t}\) (interval-censored as a category \(k_{t} \in \{1,\dots,7\}\)).
The latent 3-period history is \(H = (h_{1},h_{2},h_{3})\)
and the set of histories \(\mathcal{H}\) has 8 elements. The clock spacing
is \(\Delta = 0.25\) years. Parameters: \(\Theta = (\alpha,\theta_0,\theta_1,\pi,\varepsilon,\lambda_g,\lambda_d)\).
When we move to estimation (Section~4), we index individuals
\(i = 1,\dots,N\) with survey weights \(w_i\) and write
\(s_{it}\), \(g_{it}\), etc.\ for individual-level observations.

A \textbf{clock} (or \emph{deterministic clock}) refers to the
latent tenure or nonemployment duration that would be observed if the
true state history \(h\) were known and there were no measurement error.
Given \(h\), the clock evolves deterministically on the quarterly grid:
each continuation wave increments the clock by \(\Delta = 0.25\) years,
while a state transition resets it. For example, three consecutive waves
of true employment produce the tenure clock path
\(g^*_1, \; g^*_1 + 0.25, \; g^*_1 + 0.50\). Observed tenure may
deviate from the clock due to contamination (see below).

A \textbf{flank} is a tenure observation at one end of a latent
employment spell that spans an intermediate wave where tenure is
unobserved (because \(s_t = 0\)). For example, under state pattern
\(s = (1,0,1)\) and latent history \(h = (1,1,1)\), the spell runs
\([1,3]\) but tenure is observed only at waves~1 and~3; \(g_1\) is the
\emph{left flank} and \(g_3\) the \emph{right flank}. If both flanks
are clock-consistent, they satisfy \(g_3 = g_1 + 2\Delta\). The
spell-pair emission (Section~\ref{sec:per_spell}) exploits this
structural link to identify contamination and, indirectly,
misclassification.

A \textbf{branch} refers to one component of a mixture density.
Throughout this document, the tenure measurement model is a two-component
mixture governed by the latent contamination indicator \(m_{t}\):
the \emph{clock-consistent branch} (\(m_{t}=0\), probability
\(1-\varepsilon\)) assigns the observation to its deterministic clock
value (a point mass, i.e. a point with probability 1), while the \emph{contaminated branch}
(\(m_{t}=1\), probability \(\varepsilon\)) treats the observation as an
independent \(\mathrm{Exp}(\lambda_g)\) draw. At the spell level, the
per-wave branches combine multiplicatively across the \(K\) observed waves,
so the spell emission \eqref{eq:spell_emission} sums over all \(2^K\)
contamination patterns. Throughout this document, a wave is called
\textbf{clean} when \(m_t = 0\) (its tenure equals the clock) and
\textbf{contaminated} when \(m_t = 1\) (its tenure is an independent draw
unrelated to the clock).

\section{Tenure measurement model}

\subsection{Spell-length distribution: Exponential}

We model the latent employment-spell length as \(T_g \sim \mathrm{Exp}(\lambda_g)\)
and the latent nonemployment-spell length as \(T_d \sim \mathrm{Exp}(\lambda_d)\),
without an additive Gaussian noise component. The motivation is twofold:

\begin{itemize}
\item Identification of \((\lambda,\sigma^2)\) jointly in an EMG is delicate
  and is empirically undermined here by the point-mass at zero increment.
  Removing \(\sigma^2\) makes \(\lambda\) identification immediate from the
  empirical mean of the spell distribution.
\item The remaining tenure variation (deviations from clock-consistency)
  is captured by the contamination component \(\varepsilon\), described
  below. There is no role left for an additive Gaussian.
\end{itemize}

\subsection{Contamination model: point-mass + Exponential}

Let \(m_{t} \in \{0,1\}\) be a latent indicator for whether the wave-\(t\)
tenure report is \emph{clock-consistent}
(\(m_{t} = 0\), probability \(1-\varepsilon\)) or \emph{contaminated}
(\(m_{t} = 1\), probability \(\varepsilon\)). Contamination represents
recall error, year-rounding, dwelling-respondent substitution, or any
mechanism that produces a tenure value unrelated to the clock.

The clock-consistent branch reports the deterministic clock value exactly
(point mass). The contaminated branch reports an independent draw from the
spell-length distribution \(\mathrm{Exp}(\lambda_g)\):
\begin{equation}
g_{t} \mid h, m_{t}=0,\, s_{t}=1 \;=\; g^*_{t}(h) \quad\text{(deterministic)},
\qquad
g_{t} \mid h, m_{t}=1,\, s_{t}=1 \;\sim\; \mathrm{Exp}(\lambda_g),
\label{eq:eps_branch}
\end{equation}
where \(g^*_{t}(h)\) is the latent clock value defined below in
\eqref{eq:deterministic_clocks_eps}. This measurement model applies only
at waves where tenure is observed (\(s_{t}=1\)); when \(s_{t}=0\) the
survey records a timegap category instead and the tenure emission is not
evaluated.

The contamination indicators are independent across waves, so the joint 
probability of all the tenure observations within a spell is just the 
product of each wave's individual contribution.

\subsection{Latent clock and spell-pair structure}
\label{sec:spell_pair}

For a maximal latent E-spell spanning waves \(a \le t \le b\) (i.e.,
\(h_t = 1\) for all \(a \le t \le b\)), the latent tenure path is
\begin{equation}
g^*_a(h) = T_g, \qquad
g^*_t(h) = T_g + (t-a)\Delta \quad\text{for } a < t \le b,
\label{eq:deterministic_clocks_eps}
\end{equation}
where \(T_g \sim \mathrm{Exp}(\lambda_g)\) is a random variable giving the employment
spell-elapsed-time at wave \(a\). Under \(m_{a} = 0\), \(g_{a} = T_g\)
(observed).

A wave with \(s_{t}=0\) does not
contribute to the tenure emission (the survey records timegap, not tenure)
but its latent clock value still propagates the spell.

\subsection{Per-spell tenure emission (the key innovation)}
\label{sec:per_spell}

For each maximal latent E-spell \([a,b]\) under history \(h\), let
\(\mathcal{O} \subseteq \{a,a+1,\dots,b\}\) denote the wave indices with
\(s_{t} = 1\) (meaning \(\mathcal{O}\) gives the set of observed tenure), 
and let \(t_1 < t_2 < \cdots < t_K\) be
the elements of \(\mathcal{O}\) with \(K = |\mathcal{O}|\). Each observed
wave has a latent contamination indicator \(m_t \in \{0,1\}\), drawn
independently with \(P(m_t = 1) = \varepsilon\). Write
\(\mathbf{m} = (m_{t_1}, \dots, m_{t_K}) \in \{0,1\}^K\) for a
contamination pattern, \(\mathcal{C}(\mathbf{m}) = \{t \in \mathcal{O} :
m_t = 0\}\) for the set of \textit{clean} waves, and
\(\mathcal{D}(\mathbf{m}) = \{t \in \mathcal{O} : m_t = 1\}\) for the
\textit{contaminated} waves.

The joint tenure emission for the spell is the exact mixture over all
\(2^K\) contamination patterns:
\begin{equation}
\ell^{\mathrm{spell}}_g\!\big(\{g_t\}_{t \in \mathcal{O}} \mid h, \Theta\big)
\;=\;
\sum_{\mathbf{m} \in \{0,1\}^K}
\biggl[\prod_{t \in \mathcal{O}} \varepsilon^{m_t}(1-\varepsilon)^{1-m_t}\biggr]\,
f\!\big(\{g_t\} \mid \mathbf{m}, mathbf{h}\big),
\label{eq:spell_emission}
\end{equation}
where \(f(\{g_t\} \mid \mathbf{m}, mathbf{h})\) is the conditional density of the
observed tenures given the contamination pattern, \textbf{m}, and true employment history, 
\textbf{h}, obtained by
marginalising the latent spell-start elapsed time
\(T_g \sim \mathrm{Exp}(\lambda_g)\):

\begin{itemize}
\item If at least one wave is clean
  (\(\mathcal{C}(\mathbf{m}) \neq \emptyset\)), let
  \(t_c = \min(\mathcal{C})\) be the earliest clean wave. The clean
  report at \(t_c\) pins \(T_g = g_{t_c} - (t_c - a)\Delta\), giving
  \begin{equation}
  f \;=\; \lambda_g\, e^{-\lambda_g\bigl(g_{t_c}\, -\, (t_c-a)\Delta\bigr)}
  \cdot \prod_{\substack{t \in \mathcal{C} \\[2pt] t > t_c}}
  \!\mathcal{I}\!\big\{g_t = g_{t_c} + (t - t_c)\Delta\big\}
  \cdot \prod_{t \in \mathcal{D}} \lambda_g\, e^{-\lambda_g g_t}
  \label{eq:per_pattern_density}
  \end{equation}
  for \(g_{t_c} \ge (t_c - a)\Delta\), and zero otherwise. The first
  factor is the \(\mathrm{Exp}(\lambda_g)\) density of \(T_g\); the
  indicator product enforces clock-consistency among the remaining
  clean waves; the final product gives the independent
  \(\mathrm{Exp}(\lambda_g)\) density of each contaminated wave.
\item If all waves are contaminated
  (\(\mathcal{C}(\mathbf{m}) = \emptyset\)), \(T_g\) integrates to~1 and
  \[
  f \;=\; \prod_{t \in \mathcal{O}} \lambda_g\, e^{-\lambda_g g_t}.
  \]
\end{itemize}

\paragraph{Singletons (\(K = 1\)).}
When \(K = 1\) and \(t_1 = a\) (i.e.\ the observed wave is the spell
start), both contamination patterns yield
\(\lambda_g e^{-\lambda_g g_{t_1}}\), so \(\varepsilon\) drops out and the
emission reduces to a single \(\mathrm{Exp}(\lambda_g)\) evaluation.
When \(t_1 > a\) (the spell start is unobserved), the clean branch is
a \emph{shifted} \(\mathrm{Exp}(\lambda_g)\) density with support
\(g_{t_1} \ge (t_1-a)\Delta\), while the contaminated branch is the
unshifted \(\mathrm{Exp}(\lambda_g)\) on \([0,\infty)\); the two differ
by the factor \(e^{\lambda_g (t_1-a)\Delta}\), and \(\varepsilon\)
remains active.

\paragraph{Pairs (\(K = 2\)): the identification-critical case.}
For state pattern \(s = (1,0,1)\) under \(h = (1,1,1)\): spell \([1,3]\),
\(\mathcal{O} = \{1,3\}\), \(a = 1\), \(K = 2\). The contamination
indicators \((m_1, m_3)\) take four possible values, which are
mutually exclusive and exhaustive. By the law of total probability, the
joint emission marginalises over them: each line below is one term, prior probability of the pattern \(\times\) conditional density of
\((g_1, g_3)\) under that pattern, with the conditional density
obtained by integrating out the latent spell-start time
\(T_g \sim \mathrm{Exp}(\lambda_g)\):
\begin{align}
\ell^{\mathrm{spell}}_g \;=\;\;
  & (1-\varepsilon)^2\,\lambda_g\, e^{-\lambda_g g_1}\,
    \mathcal{I}\{g_3 = g_1 + 2\Delta\}
    && (m_1{=}0,\, m_3{=}0) \notag \\
+\; & (1-\varepsilon)\,\varepsilon\;\lambda_g^2\,
    e^{-\lambda_g(g_1 + g_3)}
    && (m_1{=}0,\, m_3{=}1) \notag \\
+\; & \varepsilon\,(1-\varepsilon)\;\lambda_g^2\,
    e^{-\lambda_g(g_1 + g_3 - 2\Delta)}
    && (m_1{=}1,\, m_3{=}0) \notag \\
+\; & \varepsilon^2\;\lambda_g^2\,
    e^{-\lambda_g(g_1 + g_3)}
    && (m_1{=}1,\, m_3{=}1)
\label{eq:spell_emission_k2}
\end{align}
where \(2\Delta = 0.5\) years. Line~1: both reports clean --- \(g_1\) pins
\(T_g\), \(g_3\) must equal \(g_1 + 2\Delta\) (point mass). Line~2: wave~1
clean (pins \(T_g = g_1\)), wave~3 an independent Exp draw. Line~3:
wave~1 contaminated, wave~3 clean --- \(T_g\) is pinned by \(g_3\):
\(T_g = g_3 - 2\Delta\), producing a \emph{shifted}
\(\mathrm{Exp}(\lambda_g)\) density for \(T_g\) (valid for
\(g_3 \ge 2\Delta\)). Line~4: both contaminated, \(T_g\) integrates out.

\paragraph{Why this is the engine of identification.}
For state pattern \(s = (1,0,1)\), the two contending latent histories
yield different \(\ell^{\mathrm{spell}}_g\):

\begin{itemize}
\item Under \(h = (1,1,1)\) (one E-spell of length~3, with miscl.\ at
  \(t=2\)): \(\mathcal{O} = \{1,3\}\), \(K = 2\). Equation
  \eqref{eq:spell_emission_k2} contributes a point mass at
  \(\{g_3 = g_1 + 2\Delta\}\) (line~1), plus three continuous terms
  (lines~2--4).
\item Under \(h = (1,0,1)\) (two separate singleton E-spells): each
  spell has \(K = 1\) with \(t_1 = a\). The emission reduces to
  \(\lambda_g e^{-\lambda_g g_t}\) for each, with no point mass and no
  dependence on \(\varepsilon\).
\end{itemize}

When \(g_3 = g_1 + 2\Delta\) (about 24.8\% of \(s=(1,0,1)\) observations
in the QLFS), the \((1,1,1)\) hypothesis fires the point mass and
dominates, pushing the posterior weight toward \((1,1,1)\) and increasing
the implied \(\pi\). When \(g_3\) is inconsistent with \(g_1 + 2\Delta\),
line~1 of \eqref{eq:spell_emission_k2} contributes zero, but the three
continuous terms (lines~2--4) remain active and still depend on
\(\varepsilon\). The assignment is then governed primarily by the Markov
prior, with secondary adjustment from the relative magnitudes of the
continuous contamination terms under \((1,1,1)\) versus the two-singleton
product \(\lambda_g^2 e^{-\lambda_g(g_1+g_3)}\) under \((1,0,1)\).

\paragraph{Triples (\(K = 3\)).}
For \(s = (1,1,1)\) under \(h = (1,1,1)\): \(\mathcal{O} = \{1,2,3\}\),
\(a = 1\), \(K = 3\). Equation~\eqref{eq:spell_emission} sums over
\(2^3 = 8\) contamination patterns following the same structure: the
all-clean pattern contributes a point mass at
\(\{g_2 = g_1 + \Delta,\; g_3 = g_1 + 2\Delta\}\); mixed patterns
contribute shifted or unshifted Exp densities depending on which wave
anchors \(T_g\); the all-contaminated pattern treats all three tenures as
independent Exp draws. This case primarily identifies \(\varepsilon\)
from the empirical fraction of \((1,1,1)\) panels whose tenure
increments are exactly \(+\Delta\) per wave.

\paragraph{Posterior contamination probability.}
For a spell with \(K \ge 2\), let
\(\tau_h^{\mathrm{spell}}\) denote the posterior probability that
\emph{at least one} of the \(K\) observed waves in the spell is
contaminated, given the observed tenures and the latent history \(h\).
By Bayes' rule, 
\begin{equation}
\tau_h^{\mathrm{spell}} \;=\;
P\!\big(\exists\, t \in \mathcal{O} : m_t = 1 \,\big|\, \{g_t\}, h, \Theta\big)
\;=\;
1 \;-\; \frac{(1-\varepsilon)^K \, f_{\mathbf{0}}}
             {\ell^{\mathrm{spell}}_g},
\label{eq:tau_spell}
\end{equation}
where \(f_{\mathbf{0}} = f(\{g_t\} \mid \mathbf{m} = \mathbf{0}, h)\) is
the all-clean conditional density (e.g. the first line of
\eqref{eq:spell_emission_k2} in the \(K=2\) case, including the
clock-consistency indicator), \((1-\varepsilon)^K\) is the prior
probability that every wave is clean, and \(\ell^{\mathrm{spell}}_g\)
from \eqref{eq:spell_emission} is the marginal likelihood summing over
all \(2^K\) contamination patterns. The numerator
\((1-\varepsilon)^K f_{\mathbf{0}}\) is therefore the joint probability
of the data and the all-clean event; dividing by
\(\ell^{\mathrm{spell}}_g\) gives the posterior of the all-clean event,
and one minus that is the posterior of its complement.

Two limits are useful. When the data exactly satisfy the clock relations
required by the all-clean branch (e.g., \(g_3 = g_1 + 2\Delta\) for
\(K=2\)), the all-clean indicator fires and \(f_{\mathbf{0}}\) is
finite, while the contaminated branches contribute only continuous
density at that point. In the limit where the point mass dominates
(formally, when the dominating measure assigns positive weight only to
the all-clean event at clock-consistent observations),
\(\tau \to 0\): the spell is judged clean with certainty. Conversely,
when the data violate every clock relation, \(f_{\mathbf{0}} = 0\) and
\(\tau = 1\): the spell must contain at least one contaminated wave.
Between these extremes \(\tau\) interpolates smoothly, and is the
quantity that \emph{aggregates} across spells in the M-step
\eqref{eq:eps_mstep} to deliver \(\hat\varepsilon\).

\subsection{N-spell timegap emission}

Nonemployment timegap is observed as a category \(k \in \{1,\dots,7\}\).
We retain the interval-censored \(\mathrm{Exp}(\lambda_d)\) emission from
the base model:
\begin{equation}
P(D \in [a_k, b_k) \mid \lambda_d) = e^{-\lambda_d a_k} - e^{-\lambda_d b_k}.
\label{eq:timegap_interval}
\end{equation}
The discrete category structure already provides robustness against small
reporting errors (since within-bin variation does not enter the likelihood).
We do not introduce an analogue of \(\varepsilon\) for timegap in this
specification. (This is an explicit design choice; see Section
\ref{sec:design_choices}.)

\subsection{Per-wave emissions for non-spell-pair cases}

For waves not falling under the per-spell emission rule (e.g., \(h_t=0\),
or \(h_t=1\) with \(s_t=0\) so tenure is unobserved), the emissions follow:

\begin{itemize}
\item \textbf{Misclassified-as-employed} (\(h_t=0,\,s_t=1\)): the observed
  \(g_t\) is drawn from the population marginal of employed tenure,
  \(g_t \sim \mathrm{Exp}(\lambda_g)\). The wave is treated as a singleton
  E-segment for emission purposes (\(K=1\), no point mass).
\item \textbf{Matched/misclassified nonemployed} (\(s_t=0\)): timegap
  emission \eqref{eq:timegap_interval} on observed \(k_{t}\).
\end{itemize}

\subsection{State misclassification} (unchanged from the base non-tenure model)

\begin{equation}
q_{ab}(\pi) = \pi^{\mathcal{I}\{a \neq b\}}\,(1-\pi)^{\mathcal{I}\{a=b\}}.
\end{equation}

\subsection{Observed-data likelihood}

For individual \(i\), let
\(y_i = \big(\{s_{it}\}_{t=1}^3,\, \{g_{it}\}_{t:s_{it}=1},\,
\{k_{it}\}_{t:s_{it}=0}\big)\) collect all observed quantities: the
three reported employment states, the reported tenures at waves where
\(s_{it}=1\), and the reported nonemployment-timegap categories
\(k_{it} \in \{1,\dots,7\}\) at waves where \(s_{it}=0\).
Given individual weight \(w_i\), the per-history conditional likelihood is
\begin{equation}
p(y_i \mid h; \Theta) =
\underbrace{\Big[\prod_{t=1}^3 q_{s_{it},h_t}(\pi)\Big]}_{\substack{\text{state}\\\text{misclassification}}} \cdot
\underbrace{\Big[\prod_{[a,b] \in \mathcal{S}_g(h)} \ell^{\mathrm{spell}}_g(\cdot)\Big]}_{\substack{\text{tenure spell}\\\text{emissions (E-spells)}}} \cdot
\underbrace{\Big[\prod_{t : h_t=0,\,s_{it}=0} P(k_{it} \mid \lambda_d)\Big]}_{\substack{\text{timegap: matched}\\\text{nonemployed}\\(h_t{=}0,\,s_t{=}0)}} \cdot
\underbrace{\Big[\prod_{t : h_t=1,\,s_{it}=0} P(k_{it} \mid \lambda_d)\Big]}_{\substack{\text{timegap: misclassified-}\\\text{as-nonemployed}\\(h_t{=}1,\,s_t{=}0)}},
\end{equation}
where \(\mathcal{S}_g(h)\) is the set of maximal E-spells in \(h\), and
\(P(k_{it} \mid \lambda_d)\) is the interval-censored
\(\mathrm{Exp}(\lambda_d)\) probability \eqref{eq:timegap_interval}
assigned to category \(k_{it}\). The first factor scores how well the
reported state sequence matches the latent history \(h\); the second
gives the joint density of all observed tenures within each E-spell of
\(h\) (using the point-mass~+~Exp mixture from \eqref{eq:spell_emission});
the third and fourth factors score the timegap reports at non-employment
waves --- separated only because the two cases generate \(k_{it}\) from
distinct latent durations (a true N-spell duration when \(h_t = 0\),
versus a fabricated/recall duration when \(h_t = 1\) but the respondent
reports nonemployment), although both currently use the same
\(\mathrm{Exp}(\lambda_d)\) emission. The
observed-data log-likelihood is
\begin{equation}
\ell(\Theta) = \sum_{i=1}^N w_i \log\bigg[\sum_{h \in \mathcal{H}} p(H_i = h \mid \alpha,\theta_0,\theta_1) \cdot p(y_i \mid h; \Theta)\bigg].
\end{equation}

\section{EM algorithm}

\subsection{Complete-data formulation}

Introduce indicators \(z_{ih} = \mathcal{I}\{H_i = h\}\) and per-spell
contamination pattern vectors \(\mathbf{m}_{i,[a,b],h} \in \{0,1\}^K\) for
each maximal E-spell of length \(K \ge 1\) under history \(h\).
The complete-data log-likelihood factors into independent blocks for
the Markov prior, state misclassification, spell contamination, and the
spell tenure / timegap emissions.

\subsection{E-step}

\paragraph{Responsibilities.}
\begin{equation}
\gamma_{ih} = \frac{p(h \mid \Theta^{\mathrm{old}}) \cdot p(y_i \mid h; \Theta^{\mathrm{old}})}{\sum_{h'} p(h' \mid \Theta^{\mathrm{old}}) \cdot p(y_i \mid h'; \Theta^{\mathrm{old}})}.
\end{equation}

\paragraph{Posterior contamination probabilities.}
For each \emph{eps-informative} E-spell under history \(h\) --- i.e.\ each
maximal E-spell \([a,b]\) with either \(K \ge 2\) observed waves, or
\(K = 1\) with \(t_1 > a\) (\(K=1\) spells with \(t_1 = a\) are not
eps-informative because both contamination patterns yield the same density)
--- the per-wave posterior contamination probability is
\begin{equation}
\nu_{it,h,[a,b]} \;=\; P(m_t = 1 \mid y_i, h, \Theta^{\mathrm{old}})
\;=\;
\frac{\sum_{\mathbf{m}:\, m_t = 1}\!
P(\mathbf{m})\, f(\{g_{it}\} \mid \mathbf{m}, h)}
{\ell^{\mathrm{spell}}_g(\{g_{it}\} \mid h, \Theta^{\mathrm{old}})},
\end{equation}
where the sum is over all contamination patterns with \(m_t = 1\), and
the denominator is the full spell emission \eqref{eq:spell_emission}.
The spell-level posterior probability that at least one wave is
contaminated is
\(\tau_{ih,[a,b]} = 1 - (1-\varepsilon)^K f_{\mathbf{0}} /
\ell^{\mathrm{spell}}_g\),
as in \eqref{eq:tau_spell}.

\subsection{M-step}

The M-step decomposes into independent blocks because the complete-data
log-likelihood is additively separable. We summarise each.

\paragraph{Markov prior block} (\(\alpha,\theta_0,\theta_1\)): unchanged
from the base free specification.
\begin{equation}
\hat\alpha = \frac{C_1}{C_0+C_1}, \qquad
\hat\theta_1 = \frac{T_{11}}{D_1}, \qquad
\hat\theta_0 = \frac{T_{01}}{D_0}.
\end{equation}

\paragraph{Misclassification} (\(\pi\)): unchanged.
\begin{equation}
\hat\pi = \frac{M}{3W}, \qquad W = \sum_i w_i.
\end{equation}

\paragraph{Tenure contamination} (\(\varepsilon\)): closed form.
The standard EM update for a per-wave Bernoulli mixture weight uses the
posterior contamination probabilities \(\nu_{it,h}\):
\begin{equation}
\hat\varepsilon \;=\;
\frac{\sum_{i,h,[a,b]:\,\mathrm{eps\text{-}inf}}\, w_i \gamma_{ih}\,
      \sum_{t \in \mathcal{O}} \nu_{it,h,[a,b]}}
     {\sum_{i,h,[a,b]:\,\mathrm{eps\text{-}inf}}\, w_i \gamma_{ih}\, K_{ih,[a,b]}}
\label{eq:eps_mstep}
\end{equation}
where the sums run over all eps-informative spells: those with \(K \ge 2\),
or \(K = 1\) with \(t_1 > a\) (the observed wave is after the spell start,
so the clean and contaminated branches differ by a shift of
\((t_1-a)\Delta\)).
\(K_{ih,[a,b]}\) is the number of waves with \(s = 1\) in the spell.
The numerator counts the posterior-expected number of contaminated
wave-tenure reports; the denominator counts all wave-tenure reports in
eps-informative spells.
Spells with \(K = 0\), or \(K = 1\) with \(t_1 = a\), contribute no
information about \(\varepsilon\) and are excluded from both sums.


\paragraph{Spell-length rate} (\(\lambda_g\)): closed form (Exponential MLE),
\begin{equation}
\hat\lambda_g = \frac{\sum_{i,h,[a,b],t\in\mathcal{O}} w_i \gamma_{ih} \cdot \xi_{it}}{\sum_{i,h,[a,b],t\in\mathcal{O}} w_i \gamma_{ih} \cdot \xi_{it} \cdot g_{it}^{\mathrm{eff}}},
\label{eq:lambda_g_mstep}
\end{equation}
where \(\xi_{it}\) is the Exp-weight assigned to observation \(t\) by the
spell-emission decomposition (1 for the leading observation of a clean
spell-pair, 1 for each observation of a contaminated spell-pair, 1 for
singletons), and \(g_{it}^{\mathrm{eff}}\) is the effective tenure (the
observed value, except for the leading wave of a clean spell-pair where
the relation \(g_t = g_a + (t-a)\Delta\) is enforced and only \(g_a\)
contributes). This is an Exponential-MLE inverse-mean.

\paragraph{Nonemployment rate} (\(\lambda_d\)): unchanged from the base EM tenure
model. Brent solver on the interval-censored Exp likelihood. The full
timegap emission machinery -- marginal interval probability
\eqref{eq:timegap_interval}, conditional transition probability for
nonemployment continuations with both flanks observed, and the start
emission for fresh nonemployment spells -- is reused verbatim from the
base / \(\rho\) implementations. This preserves apples-to-apples
comparability of \(\hat\lambda_d\) across specifications and concentrates
the structural change of Spec I in the tenure block.

\section{Identification}
\label{sec:identification_eps}

\subsection{Block separation of the complete-data log-likelihood}

The complete-data log-likelihood decomposes into four independent blocks:

\begin{enumerate}
\item Markov prior \(\rightarrow\) identifies \((\alpha,\theta_0,\theta_1)\).
\item State misclassification \(\rightarrow\) identifies \(\pi\).
\item Tenure spell-pair contamination \(\rightarrow\) identifies \(\varepsilon\).
\item Tenure spell-length and N-spell rate \(\rightarrow\) identifies
  \((\lambda_g, \lambda_d)\).
\end{enumerate}

\subsection{Identification of \(\varepsilon\)}

\(\varepsilon\) is identified from the empirical fraction of E-spell-pairs
that satisfy the exact clock-consistency relation in the data. For two
consecutive observed waves with \(h_{t-1}=h_t=1\) and \(s_{t-1}=s_t=1\),
the indicator \(\mathcal{I}\{g_t - g_{t-1} = \Delta\}\) is Bernoulli with
mean \(1-\varepsilon\). Empirically, 63.2\% of such pairs satisfy this
relation in the QLFS, giving a method-of-moments estimate \(\varepsilon \approx 0.37\).
For 2-wave bridges (\(s=(1,0,1)\) with same-spell), the indicator
\(\mathcal{I}\{g_3 - g_1 = 2\Delta\}\) has mean \((1-\varepsilon)^2\); the
empirical 24.8\% gives \(\varepsilon \approx 0.50\). The EM estimate
combines both via the responsibility-weighted average and is consistent
under standard regularity.

\subsection{Identification of \(\pi\) (the corroboration channel)}

\(\pi\) is identified from state-Markov violations weighted by tenure
flank consistency. Specifically, the M-step formula \(\hat\pi = M / (3W)\)
applies the standard counts to the responsibilities \(\gamma_{ih}\), but
\(\gamma_{ih}\) itself is now a function of the spell-pair tenure
emissions. The implications:

\begin{itemize}
\item State pattern \(s=(1,0,1)\) with \(g_3 \approx g_1 + 0.5\) yr
  (24.8\% of cases): \(\gamma\) pushed toward \(h=(1,1,1)\) by the spell-pair
  point mass. These contribute \(\pi \cdot (1-\pi)^2\) to the misclassification
  count \(M\).
\item State pattern \(s=(1,0,1)\) with \(g_3\) fresh-start-consistent
  (\(\le 3\) months, 25\% of cases): no point mass under \((1,1,1)\); both
  hypotheses yield single-Exp evaluations. \(\gamma\) split by Markov prior
  alone, with neither hypothesis dominating.
\item State pattern \(s=(0,1,0)\) with timegap-consistent flanks (e.g.,
  category 1 in both flanks): \(\gamma\) pushed toward \(h=(0,0,0)\) by the
  N-spell timegap consistency.
\end{itemize}

The result is that \(\hat\pi\) is identified from \emph{both} state-Markov
violations (the simple-model channel) \emph{and} tenure flank consistency
(the new corroboration channel). When both channels agree, \(\hat\pi\) is
sharper and slightly higher than the simple-model estimate, because tenure
flank evidence has reclassified a subset of \((s_{t-1}=0, s_t=1)\) state
transitions from genuine to misclassified.

Specification should yield \(\pi\) modestly
above the simple-model estimate of 2.5\%-3\%. Numerical simulation (a single
M-step, holding all else equal, on the QLFS sample) suggests \(\hat\pi\)
in the range 3--5\%, depending on \(\hat\varepsilon\). 
In this case, tenure data corroborates and refines the simple-model misclassification
estimate.

\subsection{Identification of \(\lambda_g\)}

\(\lambda_g\) is identified from the empirical mean of (a) wave-1 employed
tenure (\(\sim 189\)k observations), (b) misclassified-as-employed tenures
(small fraction), (c) the leading observation of clean spell-pairs, and
(d) all observations within contaminated spell-pairs. As a single-parameter
Exponential MLE, identification is immediate from the inverse-mean.

\section{Comparison with base EM tenure model}
\label{sec:design_choices}

\begin{center}
\begin{tabular}{llll}
\toprule
Feature & Base model & Contamination model \\
\midrule
Spell-length distribution & EMG\((\lambda_g,\sigma_g^2)\) & Exp\((\lambda_g)\) \\
Tenure noise & \(N(0,2\sigma_g^2)\) & point mass \(\delta(\Delta=\Delta)\) or Exp\((\lambda_g)\) through \(\varepsilon\) \\
Tenure-flank \(\to\) \(\pi\) ID & no & yes (spell-pair) \\
\(\sigma_g^2\) collapse issue & yes &  not present (no \(\sigma\)) \\
Free parameters & 7 & 7 \\
\bottomrule
\end{tabular}
\end{center}

\subsection{Why no \(\varepsilon_d\) for timegap}

The QLFS timegap variable is a 7-category interval-censored measurement.
Within-bin variation does not enter the likelihood, providing implicit
robustness against small reporting errors. The deterministic clock for
nonemployment durations advances 3 months per wave; for a 24-month timegap
at wave 1, the clock predicts category 5 ([12,36) months) at all three
waves. Bin-crossing happens primarily at the boundaries of categories
(e.g., category 1 to 2 at 3 months). Adding an \(\varepsilon_d\) parameter
would be redundant with the categorical robustness already in place and
weakens identification of \(\lambda_d\). We therefore omit it. (A robustness
check version with \(\varepsilon_d\) can be added if a referee requests it.)

\subsection{Why no \(\rho\) (full mismatch)}

In the conversation that produced this spec, the user is interested in
also reporting a \emph{unified mismatch} model where \(\rho\) governs joint
state-and-tenure mismatch (Story B). This is preserved in the existing
\(\rho\) implementation (see \texttt{EM tenure rho.tex}) and run in parallel
with Spec I. The two specifications target different stories and are
complementary:

\begin{itemize}
\item Spec I (\(\varepsilon\) only): tenure errors generate inconsistent
  reports while state is correct. \(\pi\) is corroborated and refined by
  tenure flanks.
\item \(\rho\) model: a fraction of observations is jointly random
  (state-and-tenure unrelated to true trajectory). \(\pi\) is identified
  from non-mismatched observations only.
\end{itemize}

\section{Nesting and tests}

\paragraph{Nesting.} The base specification is recovered by setting
\(\varepsilon = 0\) and replacing the spell-pair emission with an
EMG-and-Gaussian-increment evaluation. Spec I does \emph{not} formally
nest the base model because the distributional families differ
(point-mass-Exp vs.\ EMG-Gaussian); however the comparison can be made via
non-nested model selection (BIC, Vuong test).

\paragraph{LR test of \(\varepsilon = 0\).} Within Spec I, the restriction
\(\varepsilon = 0\) is testable. Under \(H_0: \varepsilon = 0\) the spell-pair
emission collapses to deterministic clock; observations with non-exact
increments contribute zero likelihood. The unrestricted Spec I dominates
trivially given the data.

\paragraph{LR test of CTMC-link.} The linked specification \(\lambda_g = -\log(\theta_1)/\Delta\)
can be applied to Spec I exactly as in the base model. The LR test is
\(\chi^2_2\) under the null.

\section{Preliminary results}
\label{sec:preliminary_results}
Here we present preliminary evidence on the model presented above. 
First, some descriprive evidence from the QLFS data are given. 
Second, the estimation of the tenure contamination model is given. 

\subsection{Descriptive evidence from the QLFS}

We now quantify the empirical regularities noted in Section~1 more
systematically. Table~\ref{tab:descriptive_tenure} reports diagnostics
computed from the three-wave QLFS panel.

\begin{table}[H]
\centering
\caption{Descriptive duration statistics from the QLFS}
\label{tab:descriptive_tenure}
\begin{tabular}{@{}lrr@{}}
\toprule
Diagnostic & $N$ & Share \\
\midrule
\multicolumn{3}{@{}l}{\textit{E$\to$E continuation increments (consecutive employed waves)}} \\[2pt]
\quad Exact $+3$-month increment ($\Delta g = \Delta$) & 343\,850 & 63.2\% \\
\quad Deviation $> 12$ months from $+\Delta$ & & 18.7\% \\[4pt]
\multicolumn{3}{@{}l}{\textit{State pattern $s = (1,0,1)$: tenure-flank classification}} \\[2pt]
\quad $g_3 \approx g_1 + 2\Delta$ & 5\,997 & 27.4\% \\
\quad $g_3 \le 3$ months (fresh-start consistent) & & 18.3\% \\[4pt]
\multicolumn{3}{@{}l}{\textit{Within-panel fresh starts ($s_{t-1}=0,\, s_t=1$)}} \\[2pt]
\quad Tenure $> 12$ months at entry wave & 35\,389 & 36.7\% \\[4pt]
\multicolumn{3}{@{}l}{\textit{State pattern $s = (0,1,0)$: timegap-flank consistency}} \\[2pt]
\quad $k_1 = k_3$ (same timegap category at both flanks) & 6\,445 & 63.7\% \\
\midrule
\multicolumn{3}{@{}l}{\textit{Method-of-moments $\hat\varepsilon$}} \\[2pt]
\quad From adjacent E$\to$E pairs: $\hat\varepsilon = 1 - 0.632$ & & 0.368 \\
\quad From $(1,0,1)$ bridges: $\hat\varepsilon = 1 - \sqrt{0.274}$ & & 0.476 \\
\bottomrule
\end{tabular}\\[4pt]
{\footnotesize\textit{Notes.} $N$ is the number of wave-pairs or panels
contributing to each diagnostic. Shares are unweighted.}
\end{table}

Three features stand out. First, 63.2\% of consecutive employed-wave pairs
report the exact clock-consistent increment of $+3$ months, while 18.7\%
deviate by more than a year. This bimodal pattern---a point mass at the
deterministic clock value plus a heavy-tailed residual---directly motivates
the point-mass + Exp contamination mixture. A simple method-of-moments
estimator gives $\hat\varepsilon \approx 0.37$ from adjacent pairs.

Second, among within-panel fresh starts ($s_{t-1}=0, s_t=1$), 36.7\%
report tenure exceeding 12~months. Since a genuine job entry cannot have
more than $\Delta = 3$~months of tenure at the survey wave, these
observations imply misclassification of employment status at $t{-}1$ and
provide a model-free lower bound on the misclassification rate among
observed $(0 \to 1)$ transitions.

Third, for the state pattern $s = (1,0,1)$, 27.4\% of panels have tenure
flanks consistent with continuous employment ($g_3 \approx g_1 + 2\Delta$),
providing direct evidence for the corroboration channel of
Section~\ref{sec:identification_eps}. An additional 18.3\% have short
tenure at wave~3 ($g_3 \le 3$~months), consistent with genuine exit and
re-entry. The remaining $\sim$54\% are ambiguous or show year-rounding
artefacts.

\subsection{Estimation results}

Table~\ref{tab:eps_estimates} reports parameter estimates from the
contamination model under specifications alongside the
baseline simple models that use employment states alone. The preferred
specification (stationary, CTMC-linked) is in the rightmost column.

\begin{table}[H]
\centering
\caption{Parameter estimates: baseline models and contamination model}
\label{tab:eps_estimates}
\footnotesize
\begin{tabular}{@{}l cc cccc@{}}
\toprule
& \multicolumn{2}{c}{Simple (no tenure)} & \multicolumn{2}{c}{Contamination model} \\
\cmidrule(lr){2-3} \cmidrule(l){4-5}
& No ME & Symm.\ ME  & Linked & \textbf{Stat--Linked} \\
\midrule
Entry rate $\theta_0$ (\%) & 8.34 & 2.90 & 4.08 & \textbf{4.14} \\
Exit rate $1{-}\theta_1$ (\%) & 8.36 & 2.91 & 4.26 & \textbf{4.26} \\
Miscl.\ rate $\pi$ (\%) & --- & 2.98  & 10.21 & \textbf{10.02} \\
Contam.\ rate $\varepsilon$ (\%) & --- & ---  & 24.79 & \textbf{24.69} \\
$\lambda_g$ & --- & --- & 0.174 & \textbf{0.174} \\
$\lambda_d$ & --- & --- & 0.167 & \textbf{0.169} \\
\addlinespace
Mean E spell (months) & --- & --- & 68.9 & \textbf{68.9} \\
Mean N spell (months) & --- & ---  & 72.0 & \textbf{71.0} \\
Init.\ empl.\ prob.\ $\alpha$ & --- & ---  & 0.550 & \textbf{0.493} \\
\addlinespace
CTMC link & --- & --- & Yes & Yes \\
Stationarity & --- & ---  & No & Yes \\
\bottomrule
\end{tabular}\\[4pt]
{\footnotesize\textit{Notes.} Simple models estimated by MLE on three-wave
employment state sequences ($N=22{,}779$). Contamination models estimated
by EM on three-wave panels with tenure and timegap ($N=402{,}923$).
``Linked'' constrains $\lambda_g = -\log\theta_1/\Delta$ and
$\lambda_d = -\log(1-\theta_0)/\Delta$. ``Stationary'' constrains
$\alpha = \theta_0/(\theta_0 + 1 - \theta_1)$.}
\end{table}

\paragraph{Contamination rate.}
$\hat\varepsilon$ is stable across all variants at
$\approx 24$--$25\%$, below the method-of-moments estimates (36.8\%
adjacent, 47.6\% bridges) in Table~\ref{tab:descriptive_tenure}. The
gap arises because the MoM attributes \emph{all} non-exact increments
to contamination, whereas the EM accounts for the full mixture of latent
histories: some non-exact increments arise from genuine transitions (new
spells), not contamination within a continuing spell.

\paragraph{Misclassification rate.}
The contamination model yields $\hat\pi \approx 10\%$ under the linked
specifications, substantially above the simple ME model's 2.98\%. This
reflects the corroboration channel: tenure-flank consistency in the state
pattern $s = (1,0,1)$ provides direct evidence that some observed
$(1 \to 0 \to 1)$ sequences are misclassified continuous employment. The
simple model, lacking tenure information, identifies $\pi$ from the
$2^3 = 8$ state-pattern frequencies alone. The descriptive finding that
36.7\% of within-panel fresh starts have impossibly long tenure
(Table~\ref{tab:descriptive_tenure}) is consistent with a
misclassification rate well above the simple model's 3\%.

\paragraph{Transition rates.}
Under the preferred stationary-linked specification, the entry rate is
$\hat\theta_0 = 4.14\%$ and the exit rate is
$1 - \hat\theta_1 = 4.26\%$, yielding a steady-state employment rate of
$\hat\alpha = 49.3\%$. These fall between the na\"ive no-ME rates
($\sim$8.3\%) and the simple ME rates ($\sim$2.9\%). The na\"ive model
overstates transitions because observed state changes include
misclassifications. The simple ME model understates them because, without
tenure information, it assigns too little probability to misclassification
and consequently deflates implied true transitions.


\section{Summary}

The contamination model replaces the Gaussian-EMG measurement model with a structurally
correct point-mass + Exp contamination mixture and introduces a spell-pair
joint emission that allows tenure flanks to identify misclassification.
The net result is no
\(\sigma_g^2\) collapse issue, and a direct identification channel for
\(\pi\) via tenure flank consistency in same-spell wave pairs.

\end{document}
